\chapter{Практика}

\section{Задача}
Определить последовательность действий устройства при отправке пакета устройству из своей сети и из внешней сети.

\section{Принятие решения о том, в какую сеть будет
отправлен пакет}

\subsection{Устройства в одной сети}

\begin{figure}[H] 
    \centering
    \includegraphics[width=0.75\textwidth]{nw2.png}
    \caption{Топология сети}
\end{figure}

Допустим, отправка идет от PC1 к PC2. \\

Логика работы PC1:

\begin{enumerate}
    \item PC1 берет свой IP и маску сети, определяет с помощью побитового "И", к какой сети он принадлежит.
    \item PC1 берет IP получателя и маску сети, определяет с помощью побитового "И", к какой сети принадлежит получатель.
    \item PC1 побитово сравнивает результаты с пункта 1 и 2, видит совпадение и понимает, что получатель в той же сети.
    \item PC1 смотрит ARP кеш с надеждой найти нужный MAC адрес по IP. Если он его находит, то отправляет данные на этот MAC, в ином случае с помощью ARP ищет нужный MAC.  
\end{enumerate}

\subsection{Устройства в разных сетях}

\begin{figure}[H] 
    \centering
    \includegraphics[width=0.75\textwidth]{nw3.png}
    \caption{Топология сети}
\end{figure}

Допустим, отправка идет от PC1 к PC3. \\

Логика работы PC1, R1, R2:

\begin{enumerate}
    \item PC1 берет свой IP и маску сети, определяет с помощью побитового "И", к какой сети он принадлежит.
    \item PC1 берет IP получателя и маску сети, определяет с помощью побитового "И", к какой сети принадлежит получатель.
    \item PC1 побитово сравнивает результаты с пункта 1 и 2, видит несовпадение и понимает, что получатель в другой сети.
    \item PC1 отправляет данные на R1.
    \item R1 принимает данные, "распаковывает" их до L3 уровня, смотрит IP адрес назначения, выполняет побитовое "И" с масками всех подключенных сетей.
    \item R1 не находит нужную сеть и идет в свою таблицу маршрутизации и ищет, как попасть в сеть назначения и понимает, что сделать это можно через R2.
    \item R1 собирает снова кадр и отправляет на R2.
    \item R2 принимает данные, "распаковывает" их до L3 уровня, смотрит IP адрес назначения, выполняет побитовое "И" с масками всех подключенных сетей.
    \item R2 находит нужную сеть и идет в свою таблицу маршрутизации и узнает, что устройство подключено напрямую, и отправляет данные на нужный интерфейс.
    \item R2 собирает кадр и отправляет на PC3.
\end{enumerate}

\endinput